\chapter{Dasar Teori}

\section{Konsep Dasar Otomatisasi Perangkat Lunak}

\subsection{\textit{Continuous Integration} (CI) dan \textit{Continuous Deployment} (CD) pada Ekosistem Frontend}
\textit{Continuous Integration} (CI) dan \textit{Continuous Deployment} (CD) adalah metodologi krusial dalam \textit{DevOps} untuk mengotomatisasi pengujian dan peluncuran kode. Berbeda dengan \textit{backend} yang membutuhkan alokasi \textit{server} yang memproses logika bisnis secara terus-menerus, \textit{pipeline} CI/CD \textit{frontend} dirancang untuk melakukan kompilasi dari kode sumber menjadi sekumpulan aset statis (HTML, CSS, JavaScript). Alur ini umumnya melewati "gerbang kualitas" (\textit{quality gates}) seperti \textit{linting}, \textit{unit testing}, \textit{build}, dan akhirnya didistribusikan ke jaringan CDN. Otomatisasi ini secara empiris mempercepat iterasi pengembangan dan mengurangi insiden \textit{bug} di ranah produksi [\ref{cicd}].

\begin{figure}[H]
    \centering
    \includegraphics[width=0.9\textwidth]{cicd-diagram.png}
    \caption{Ilustrasi alur Continuous Integration dan Continuous Deployment}
\end{figure}

\clearpage
\subsection{Framework Vue.js}
Vue.js adalah \textit{progressive JavaScript framework} yang ideal untuk membangun antarmuka web interaktif yang modern. Berdasarkan kajian literatur, Vue memiliki keunggulan arsitektural yang memungkinkan pemisahan antara tampilan (UI) dan logika data secara deklaratif, dengan kurva pembelajaran yang lebih landai dibanding kerangka kerja lainnya [\ref{vue}]. Pada lingkungan pengembangan modern, Vue dikompilasi menggunakan ekosistem \textit{build tool} (seperti Vite) untuk menghasilkan aset distribusi statis.

\begin{figure}[H]
    \centering
    \includegraphics[width=0.7\textwidth]{vue.png}
    \caption{Arsitektur komponen dan reaktivitas data pada ekosistem Vue.js}
\end{figure}

\clearpage
\section{Infrastruktur dan Lingkungan Deployment}

\subsection{Otomatisasi dengan GitHub Actions}
GitHub Actions menyediakan infrastruktur \textit{pipeline} CI/CD yang terintegrasi penuh (\textit{native}) di dalam ekosistem GitHub. Sistem ini mengeksekusi instruksi dari berkas berformat YAML (\texttt{.yml}) di dalam wadah virtual (\textit{runner}) berbasis \textit{event}, seperti pembaruan kode (\textit{commit}) atau pembuatan permintaan tarik (\textit{Pull Request}).

\begin{figure}[H]
    \centering
    \includegraphics[width=0.7\textwidth]{github-actions.png}
    \caption{Konsep dasar dan siklus eksekusi workflow pada GitHub Actions}
\end{figure}

\subsection{Deployment dengan Vercel}
Vercel adalah platform pengiriman berbasis komputasi awan yang dioptimalkan untuk menge-host aplikasi \textit{frontend}. Platform ini menyediakan sistem \textit{Content Delivery Network} (CDN) yang memungkinkan aset statis hasil kompilasi aplikasi untuk diakses secara global dengan kecepatan maksimal dan waktu henti (\textit{downtime}) minimal.

\begin{figure}[H]
    \centering
    \includegraphics[width=0.8\textwidth]{vercel.png}
    \caption{Distribusi aset statis frontend secara global menggunakan infrastruktur CDN Vercel}
\end{figure}


\chapter{Hasil dan Pembahasan}

\section{Persiapan Lingkungan dan Konfigurasi Deployment}

\subsection{Konfigurasi Environment Deployment \& GitHub Secrets}
Langkah awal dalam perancangan \textit{deployment} Vercel tanpa interaksi antarmuka otomatis (\textit{headless}) adalah menyiapkan file \texttt{vercel.json}. File ini berfungsi mengatur \textit{routing} dasar pada server Vercel. Untuk memberikan otorisasi kepada GitHub Actions agar dapat mendeploy kode secara langsung, sebuah kunci otentikasi yaitu \texttt{VERCEL\_TOKEN} di-\textit{generate} secara manual pada pengaturan akun Vercel. Kunci ini merupakan kredensial sangat rahasia yang tidak boleh terekspos ke publik.

\begin{figure}[H]
    \centering
    \includegraphics[width=0.9\textwidth]{1.png}
    \caption{Konfigurasi file vercel.json untuk routing Vercel}
    
    \vspace{0.5cm}
    
    \includegraphics[width=0.9\textwidth]{10.png}
    \caption{Pembuatan kredensial VERCEL\_TOKEN di dashboard akun Vercel}
\end{figure}

Agar proses \textit{deploy} mengarah ke entitas proyek yang tepat, diperlukan identifikasi unik berupa \texttt{VERCEL\_ORG\_ID} dan \texttt{VERCEL\_PROJECT\_ID}. Kedua nilai ini diekstraksi melalui antarmuka baris perintah dengan mengeksekusi perintah \texttt{npx vercel link}. Setelah semua variabel didapatkan, mereka disematkan dengan aman ke dalam \textit{Repository Secrets} di GitHub. Melalui konfigurasi rahasia ini, otomatisasi penyebaran dikendalikan murni oleh skrip \textit{pipeline} GitHub Actions tanpa menghubungkan repositori Git secara langsung ke \textit{dashboard} Vercel.

\begin{figure}[H]
    \centering
    \includegraphics[width=0.9\textwidth]{11.png}
    \caption{Pengekstraksian VERCEL\_ORG\_ID dan VERCEL\_PROJECT\_ID menggunakan CLI}
\end{figure}

\begin{figure}[H]
    \centering
    \includegraphics[width=0.9\textwidth]{12.png}
    \caption{Akses menu Repository Secrets pada pengaturan proyek di GitHub}
\end{figure}

\begin{figure}[H]
    \centering
    \includegraphics[width=0.9\textwidth]{13.png}
    \caption{Penyematan tiga variabel rahasia Vercel di GitHub Secrets}
\end{figure}

\begin{figure}[H]
    \centering
    \includegraphics[width=0.9\textwidth]{14.png}
    \caption{Tampilan proyek pada dashboard Vercel sebelum terintegrasi penuh}
\end{figure}

\clearpage
\section{Pengembangan Vue.js dan Manajemen Pipeline}

\subsection{Struktur Proyek Vue.js dan Pengujian Lokal}
Sebelum diintegrasikan ke dalam ekosistem CI/CD, arsitektur dasar dari \textit{frontend} Vue.js disusun dengan pemisahan direktori komponen dan rute secara rapi. Pengujian lokal dilakukan menggunakan fitur peramban lokal (\textit{localhost}) untuk memverifikasi bahwa antarmuka halaman utama (Homeview) ter-render secara utuh tanpa ada dependensi komponen yang rusak.

\begin{figure}[H]
    \centering
    \includegraphics[width=0.55\textwidth]{6.png}
    \caption{Struktur hierarki folder dan file proyek frontend Vue.js}
\end{figure}

\begin{figure}[H]
    \centering
    \includegraphics[width=0.9\textwidth]{7.png}
    \caption{Tampilan render halaman Homeview saat diuji coba pada localhost}
\end{figure}

Bagian integral dari pengujian fungsionalitas \textit{frontend} adalah koneksinya terhadap antarmuka pemrograman aplikasi (API) dari \textit{backend} Laravel. Pada kondisi ideal, halaman Portofolio sukses mengambil dan merender data dari Laravel. Namun, aplikasi ini juga telah mengimplementasikan strategi penanganan galat (\textit{error handling}) atau \textit{graceful degradation}. Ketika \textit{server} Laravel dimatikan atau terputus, aplikasi Vue tidak mengalami \textit{crash} aplikasi layar putih, melainkan mampu menangkap status galat tersebut dan menampilkannya sebagai umpan balik visual yang jelas bagi pengguna.

\begin{figure}[H]
    \centering
    \includegraphics[width=0.9\textwidth]{8.png}
    \caption{Pengujian kelancaran komunikasi API dengan backend Laravel pada halaman Portofolio}
\end{figure}

\begin{figure}[H]
    \centering
    \includegraphics[width=0.9\textwidth]{9.png}
    \caption{Tampilan penanganan galat (error handling) saat backend Laravel sengaja diputus}
\end{figure}

\subsection{Perancangan Pipeline CI/CD (frontend.yml)}
Pusat dari otomatisasi pengujian didefinisikan dalam file \texttt{frontend.yml}. Arsitektur skrip ini dibangun menggunakan empat buah "gerbang kualitas" (\textit{quality gates}) utama yang dieksekusi secara berurutan. Tahap pertama, \texttt{lint}, memvalidasi standardisasi sintaksis kode. Tahap kedua, \texttt{test}, menggunakan peranti pengujian Vitest untuk mengeksekusi skenario uji unit; mekanisme \textit{fail-fast} berlaku di mana kegagalan pengujian akan menghentikan seluruh \textit{pipeline}. 

\begin{figure}[H]
    \centering
    \includegraphics[width=0.5\textwidth]{2.png}
    \caption{Skrip konfigurasi job lint untuk validasi sintaks}
\end{figure}

\begin{figure}[H]
    \centering
    \includegraphics[width=0.5\textwidth]{3.png}
    \caption{Skrip konfigurasi job test menggunakan peranti Vitest}
\end{figure}

Tahap ketiga, \texttt{build}, mengompilasi kode Vue menjadi berkas siap rilis di folder \texttt{dist/}. Terakhir, tahap \texttt{deploy} memanfaatkan kredensial \textit{Secrets} dan antarmuka CLI Vercel untuk meluncurkan hasil kompilasi \texttt{dist/} ke \textit{server} publik, dengan batasan ketat bahwa proses ini hanya terpicu jika perubahan masuk ke dalam \textit{branch} utama (\texttt{main}).

\begin{figure}[H]
    \centering
    \includegraphics[width=0.6\textwidth]{4.png}
    \caption{Skrip konfigurasi job build untuk kompilasi aset frontend}
\end{figure}

\begin{figure}[H]
    \centering
    \includegraphics[width=0.6\textwidth]{5.png}
    \caption{Skrip konfigurasi job deploy ke platform Vercel khusus untuk branch main}
\end{figure}

\clearpage
\section{Eksekusi Pipeline dan Kolaborasi Repositori}

\subsection{Simulasi Fail-Fast dan Penanganan Flagging Repositori}
Implementasi \textit{pipeline} dalam skenario dunia nyata menemui kendala khusus; repositori resmi organisasi mata kuliah (\texttt{KEPL2026}) sedang berstatus \textit{flag} (peninjauan), yang mengakibatkan GitHub menonaktifkan fitur Actions secara paksa. Oleh karena itu, uji coba eksekusi dialihkan seluruhnya ke dalam repositori mandiri pada organisasi pribadi \texttt{AcademyCity-RND}. Pengujian reliabilitas dimulai pada \textit{branch} \texttt{feat} dengan mensimulasikan kegagalan \textit{unit test} secara sengaja. Log \textit{runner} menangkap kegagalan asersi Vitest yang menyebabkan \textit{pipeline} berwarna merah dan memblokir (\textit{skip}) tahapan \textit{Build} dan \textit{Deploy}. Setelah kesalahan logika diatasi, sistem CI memberikan verifikasi hijau yang berarti kode lolos uji kualitas.

\begin{figure}[H]
    \centering
    \includegraphics[width=0.9\textwidth]{15.png}
    \caption{Simulasi pipeline gagal (fail-fast) akibat malfungsi unit test pada branch feat}
\end{figure}

\begin{figure}[H]
    \centering
    \includegraphics[width=0.6\textwidth]{16.png}
    \caption{Pemeriksaan detail log runner yang menunjukkan kegagalan asersi pada Vitest}
\end{figure}

\begin{figure}[H]
    \centering
    \includegraphics[width=0.6\textwidth]{17.png}
    \caption{Verifikasi hijau pada pipeline pasca perbaikan sintaks kode}
\end{figure}

\subsection{Manajemen Pull Request dan Environment Protection}
Alur kolaborasi dilanjutkan dengan pembuatan \textit{Pull Request} bertahap. Kode dipromosikan dari \textit{branch} \texttt{feat} menuju \texttt{dev}, di mana \textit{pipeline} CI sukses menguji kode, namun mengabaikan proses \textit{deploy} karena adanya aturan proteksi cabang. Promosi tahap akhir dilakukan melalui \textit{Pull Request} dari \texttt{dev} menuju \texttt{main}. Pada tahap produksi ini, diaktifkan pengamanan ekstra berupa kewajiban persetujuan (\textit{approval}) manual dari seorang peninjau (\textit{reviewer}), mencegah masuknya kode ke lingkungan publik secara otomatis tanpa pengawasan.

\begin{figure}[H]
    \centering
    \includegraphics[width=0.9\textwidth]{18.png}
    \caption{Proses pembuatan Pull Request dari branch feat menuju dev}
\end{figure}

\begin{figure}[H]
    \centering
    \includegraphics[width=0.9\textwidth]{19.png}
    \caption{Pengujian pipeline berjalan di dev dengan status deploy tertahan (skip)}
\end{figure}

\begin{figure}[H]
    \centering
    \includegraphics[width=0.9\textwidth]{20.png}
    \caption{Inisiasi pembuatan Pull Request final dari branch dev menuju main}
\end{figure}

\begin{figure}[H]
    \centering
    \includegraphics[width=0.9\textwidth]{21.png}
    \caption{Penahanan pipeline pada tahap deploy akibat kebutuhan Required Reviewers}
\end{figure}

\begin{figure}[H]
    \centering
    \includegraphics[width=0.9\textwidth]{22.png}
    \caption{Proses pemberian persetujuan (approval) manual oleh reviewer yang berwenang}
\end{figure}

\subsection{Analisis Log Eksekusi Sukses}
Setelah persetujuan (\textit{approval}) diberikan, seluruh siklus \textit{job} dieksekusi oleh GitHub Actions tanpa hambatan. Log eksekusi terminal untuk masing-masing tahap—mulai dari proses instalasi dependensi, eksekusi perintah \textit{lint} dan \textit{test}, pembangunan (\textit{build}) arsitektur, hingga sinkronisasi data statis dengan Vercel—menunjukkan status keberhasilan penuh (\textit{success}).

\begin{figure}[H]
    \centering
    \includegraphics[width=0.9\textwidth]{23.png}
    \caption{Keseluruhan gerbang kualitas (quality gates) berhasil terlewati di branch main}
\end{figure}

\begin{figure}[H]
    \centering
    \includegraphics[width=0.8\textwidth]{24.png}
    \caption{Log terminal memperlihatkan kelancaran proses Linting}
\end{figure}

\begin{figure}[H]
    \centering
    \includegraphics[width=0.8\textwidth]{25.png}
    \caption{Log terminal membuktikan kesuksesan proses Unit Testing}
\end{figure}

\begin{figure}[H]
    \centering
    \includegraphics[width=0.8\textwidth]{26.png}
    \caption{Log terminal menunjukkan tahap Build yang berhasil mengompilasi aset statis}
\end{figure}

\begin{figure}[H]
    \centering
    \includegraphics[width=0.8\textwidth]{27.png}
    \caption{Log terminal mengonfirmasi peluncuran aset sukses ke infrastruktur Vercel}
\end{figure}

\section{Hasil Akhir Penyebaran (Deployment)}

\subsection{Verifikasi Live Deployment pada Infrastruktur Cloud}
Verifikasi pasca-peluncuran membandingkan panel instrumen Vercel pada dua lingkungan organisasi. Terlihat repositori pribadi berhasil menautkan distribusi statisnya, bertolak belakang dengan lingkungan repositori organisasi resmi yang masih vakum akibat masalah administratif (\textit{flagging}). Saat tautan publik diakses langsung via peramban, \textit{frontend} beroperasi secara optimal dan sangat responsif. Verifikasi lebih lanjut pada halaman Portofolio mendemonstrasikan bahwa sistem penanganan galat berfungsi persis sama seperti pada pengujian lokal; antarmuka dapat merespons putusnya koneksi dari sisi \textit{backend} Laravel dengan stabil, memunculkan notifikasi ketiadaan data tanpa mengacaukan struktur DOM (\textit{Document Object Model}) halaman.

\begin{figure}[H]
    \centering
    \includegraphics[width=0.8\textwidth]{28.png}
    \caption{Komparasi status proyek Vercel pada repositori organisasi pribadi (kiri) versus resmi (kanan)}
\end{figure}

\begin{figure}[H]
    \centering
    \includegraphics[width=0.8\textwidth]{29.png}
    \caption{Demonstrasi halaman Homeview yang berhasil live dan diakses publik via Vercel}
\end{figure}

\begin{figure}[H]
    \centering
    \includegraphics[width=0.8\textwidth]{30.png}
    \caption{Tampilan halaman Portofolio secara aman menangani absennya pasokan data dari Laravel}
\end{figure}


\chapter{Kesimpulan}
Berdasarkan praktikum mengenai perancangan dan implementasi \textit{Continuous Deployment} (CD) pada \textit{frontend} Vue.js yang telah dilakukan, dapat ditarik beberapa kesimpulan utama sebagai berikut:
\begin{enumerate}
    \item Pemisahan konfigurasi lingkungan menggunakan \textit{GitHub Secrets} dan pengaturan \texttt{vercel.json} berhasil mengisolasi kredensial proyek dari potensi kebocoran publik. Teknik ini berhasil memfasilitasi alur \textit{deployment headless} yang sepenuhnya digerakkan oleh \textit{runner} GitHub Actions, meniadakan kebutuhan tautan repositori secara langsung ke penyedia layanan komputasi awan.
    \item Pemanfaatan arsitektur otomasi CI/CD berbasis \texttt{frontend.yml} yang terdiri atas empat \textit{quality gates} (yakni \textit{lint}, \textit{test}, \textit{build}, dan \textit{deploy}) terbukti sangat andal dalam mendeteksi dan menolak intervensi kode yang cacat secara \textit{fail-fast}, menghemat sumber daya peladen secara signifikan.
    \item Konfigurasi keamanan berlapis melalui \textit{Environment Protection} pada cabang \texttt{main} mendemonstrasikan kelayakan praktik terbaik dalam \textit{DevSecOps}, dengan memberlakukan kontrol rilis manual (\textit{approval}) terakhir sebelum aset aplikasi dipublikasikan ke publik.
    \item Desain kode Vue.js yang tangguh dan bersifat toleran terhadap kesalahan eksternal fungsionalitas API \textit{backend} (\textit{graceful degradation}) memegang peranan vital dalam memastikan bahwa keterbatasan infrastruktur tidak berdampak pada \textit{crash}-nya aplikasi di sisi klien.
    \item Strategi tanggap darurat yang mengalihkan sementara pusat eksekusi \textit{pipeline} menuju repositori lingkungan pribadi (\texttt{AcademyCity-RND}) akibat insiden \textit{flagging} menunjukkan sifat ketangguhan (\textit{resilience}) dari kontrol versi Git terdistribusi.
\end{enumerate}


\chapter{Daftar Pustaka}

\begin{enumerate}
    \item \label{cicd} R. T. P. et al. \textit{Empirical analysis of GitHub Actions workflows in modern software development}. ResearchGate, 2023. [Online]. Tersedia: \url{https://www.researchgate.net/search.Search.html?query=Empirical+analysis+of+GitHub+Actions+workflows+in+modern+software+development}
    
    \item \label{vue} A. S. \& J. M. \textit{Comparative Analysis of Progressive JavaScript Frameworks: Vue.js vs React}. ResearchGate, 2022. [Online]. Tersedia: \url{https://www.researchgate.net/search.Search.html?query=Comparative+Analysis+of+Progressive+JavaScript+Frameworks:+Vue.js+vs+React}
    
    \item \label{sourcecode} Source Code Praktikum:
    \begin{itemize}
        \item Repo Matkul KEPL2026:\\
        \url{https://github.com/KEPL2026/evolusi-pl-24-542159-SV-25009-v1}
        
        \item Repo Pribadi:\\
        \url{https://github.com/AcademyCity-RND/evolusi-pl-24-542159-SV-25009-v2}
    \end{itemize}
\end{enumerate}
